\section{张量积的乘积与直和}

记号约定:$A$是环,$\left(E_{i}\right)_{i\in I},\left(E_{i}^{\prime}\right)_{i\in I}$是两族右$A$-模,$\left(F_{j}\right)_{j\in J},\left(F_{j}^{\prime}\right)_{j\in J}$是两族左$A$-模,$E$是右$A$-模,$F$是左$A$-模.

\begin{proposition}{模的乘积的张量积映射}{}
	$\left(\prod\limits_{i\in I}E_{i}\right)\otimes_{A}\left(\prod\limits_{j\in J}F_{j}\right)\to\prod\limits_{i\in I,j\in J}\left(E_{i}\otimes_{A}F_{j}\right),\left(x_{i}\right)_{i\in I}\otimes_{A}\left(y_{j}\right)_{j\in J}\mapsto\left(x_{i}\otimes_{A}y_{j}\right)_{i\in I,j\in J}$是典范的$\Z$-线性映射.
\end{proposition}

\begin{proposition}{模的直和的张量积映射}{}
	$\left(\bigboxplus\limits_{i\in I}E_{i}\right)\otimes_{A}\left(\bigboxplus\limits_{j\in J}F_{j}\right)\to\bigboxplus\limits_{i\in I,j\in J}\left(E_{i}\otimes_{A}F_{j}\right),\left(x_{i}\right)_{i\in I}\otimes_{A}\left(y_{j}\right)_{j\in J}\mapsto\left(x_{i}\otimes_{A}y_{j}\right)_{i\in I,j\in J}$是典范的$\Z$-线性映射.
\end{proposition}

\begin{proposition}{直和的张量积的函子性}{}
	设$\forall i\in I\left(f_{i}\in\Hom_{\mathsf{Eod}-A}\left(E_{i},E_{i}^{\prime}\right)\right),\forall j\in J\left(g_{j}\in\Hom_{A-\mathsf{Eod}}\left(F_{j},F_{j}^{\prime}\right)\right)$,则下图交换(其中$\phi,\phi^{\prime}$是典范映射).
	\begin{center}
		\begin{tikzcd}
			{\bigboxplus\limits_{i\in I}E_{i}\otimes_{A}\bigboxplus\limits_{j\in J}F_{j}} && {\bigboxplus\limits_{i\in I,j\in J}\left(E_{i}\otimes_{A}F_{j}\right)} \\
			\\
			{\bigboxplus\limits_{i\in I}E_{i}^{\prime}\otimes_{A}\bigboxplus\limits_{j\in J}F_{j}^{\prime}} && {\bigboxplus\limits_{i\in I,j\in J}\left(E_{i}^{\prime}\otimes_{A}F_{j}^{\prime}\right)}
			\arrow["\phi", from=1-1, to=1-3]
			\arrow["{\left(\bigboxplus\limits_{i\in I}f_{i}\right)\otimes\left(\bigboxplus\limits_{j\in J}g_{j}\right)}"', from=1-1, to=3-1]
			\arrow["{\bigboxplus\limits_{i\in I,j\in J}\left(f_{i}\otimes g_{j}\right)}", from=1-3, to=3-3]
			\arrow["{\phi^{\prime}}"', from=3-1, to=3-3]
		\end{tikzcd}
	\end{center}
\end{proposition}

\begin{corollary}{自由模的张量积的单侧展开}{}
	若$F$有一个基$\left(b_{i}\right)_{i\in I}$,则$\left(b_{i}\right)_{i\in I}$确定的典范映射$E\otimes_{A}F\to E^{\left(I\right)}$是$\Z$-模同构.
\end{corollary}

\begin{proposition}{自由模的张量积的双侧展开}{}
	设$A$是交换环,$E,F$各自的基是$\left(a_{i}\right)_{i\in I},\left(b_{j}\right)_{j\in J}$,则$\left(a_{i}\otimes_{A}b_{j}\right)_{i\in I,j\in J}$是$E\otimes_{A}F$的基.
\end{proposition}

\begin{remark}{}{}
	记号同上.我们通常滥用术语,称$\left(a_{i}\otimes_{A}b_{j}\right)_{i\in I,j\in J}$是$\left(a_{i}\right)_{i\in I}$和$\left(b_{j}\right)_{j\in J}$的张量积.
\end{remark}

\begin{corollary}{}{}
	设$\phi:\left(\prod\limits_{i\in I}E_{i}\right)\otimes_{A}F\to\prod\limits_{i\in I}\left(E_{i}\otimes_{A}F\right)$是典范映射.若$F$自由(相对的,有限生成且自由),则$\phi$是单射(相对的,双射).
\end{corollary}

\begin{corollary}{}{}
	设$A$是整环,$E,F$是自由模,则$\forall x\in E\forall y\in F\left(x\otimes_{A}y=0\implies\left(x=0\right)\vee\left(y=0\right)\right)$.
\end{corollary}

\begin{corollary}{直因子的张量积}{}
	设$M,N$分别是$E,F$的直因子,$\phi:M\otimes_{A}N\to E\otimes_{A}F$是典范映射,则$\phi$是单射,$\im\left(\phi\right)$是$E\otimes_{A}F$的直因子.
\end{corollary}

\begin{corollary}{投射摸保持单性}{}
	设$P$是投射左$A$-模,$F_{1},F_{2}$是右$A$-模,则对诸$f\in\Hom_{\mathsf{Mod}-A}\left(F_{1},F_{2}\right)$,$f$是单射$\implies$ $f\otimes\id_{P}$是单射.
\end{corollary}

\begin{corollary}{张量积保持投射摸}{}
	设$A$是交换环,$E$和$F$是投射模,则$E\otimes_{A}F$是投射模.
\end{corollary}

\begin{remark}{}{}
	设对诸$i\in I$,$E_{i}^{\prime}$是$E_{i}$的子模,对任一$j\in J$,$F_{j}^{\prime}$是$F_{j}$的子模.考虑如下典范映射:
	\begin{itemize}
		\item $\iota:\left(\bigboxplus\limits_{i\in I}E_{i}^{\prime}\right)\otimes_{A}\left(\bigboxplus\limits_{j\in J}F_{j}^{\prime}\right)\to\bigboxplus\limits_{i\in I}E_{i}\otimes_{A}\bigboxplus\limits_{j\in J}F_{j}$.
		\item 对诸$i\in I,j\in J$,记$\iota_{i,j}:E_{i}^{\prime}\otimes_{A}F_{j}^{\prime}\to E_{i}\otimes_{A}F_{j}$.
	\end{itemize}
	那么,我们有典范的$\Z$-模同构$\im\left(\iota\right)\simeq\bigboxplus\limits_{i\in I,j\in J}\im\left(\iota_{i,j}\right)$.
\end{remark}
